\documentclass[a4paper,12pt,twoside]{memoir}

% Castellano
\usepackage[spanish,es-tabla,es-nodecimaldot]{babel}
\selectlanguage{spanish}
\usepackage[utf8]{inputenc}
\usepackage[T1]{fontenc}
\usepackage{lmodern} % Scalable font
\usepackage{microtype}
\usepackage{placeins}
\usepackage{enumitem}

\RequirePackage{booktabs}
\RequirePackage[table]{xcolor}
\RequirePackage{xtab}
\RequirePackage{multirow}

% Links
\PassOptionsToPackage{hyphens}{url}\usepackage[colorlinks]{hyperref}
\hypersetup{
	allcolors = {red}
}

% Ecuaciones
\usepackage{amsmath}

% Rutas de fichero / paquete
\newcommand{\ruta}[1]{{\sffamily #1}}

% Párrafos
\nonzeroparskip

% Huérfanas y viudas
\widowpenalty100000
\clubpenalty100000

% Imágenes

% Comando para insertar una imagen en un lugar concreto.
% Los parámetros son:
% 1 --> Ruta absoluta/relativa de la figura
% 2 --> Texto a pie de figura
% 3 --> Tamaño en tanto por uno relativo al ancho de página
\usepackage{graphicx}
\usepackage{subcaption}
\newcommand{\imagen}[3]{
	\begin{figure}[!h]
		\centering
		\includegraphics[width=#3\textwidth]{#1}
		\caption{#2}\label{fig:#1}
	\end{figure}
	\FloatBarrier
}

% Comando para insertar una imagen sin posición.
% Los parámetros son:
% 1 --> Ruta absoluta/relativa de la figura
% 2 --> Texto a pie de figura
% 3 --> Tamaño en tanto por uno relativo al ancho de página
\newcommand{\imagenflotante}[3]{
	\begin{figure}
		\centering
		\includegraphics[width=#3\textwidth]{#1}
		\caption{#2}\label{fig:#1}
	\end{figure}
}

% El comando \figura nos permite insertar figuras comodamente, y utilizando
% siempre el mismo formato. Los parametros son:
% 1 --> Porcentaje del ancho de página que ocupará la figura (de 0 a 1)
% 2 --> Fichero de la imagen
% 3 --> Texto a pie de imagen
% 4 --> Etiqueta (label) para referencias
% 5 --> Opciones que queramos pasarle al \includegraphics
% 6 --> Opciones de posicionamiento a pasarle a \begin{figure}
\newcommand{\figuraConPosicion}[6]{%
  \setlength{\anchoFloat}{#1\textwidth}%
  \addtolength{\anchoFloat}{-4\fboxsep}%
  \setlength{\anchoFigura}{\anchoFloat}%
  \begin{figure}[#6]
    \begin{center}%
      \Ovalbox{%
        \begin{minipage}{\anchoFloat}%
          \begin{center}%
            \includegraphics[width=\anchoFigura,#5]{#2}%
            \caption{#3}%
            \label{#4}%
          \end{center}%
        \end{minipage}
      }%
    \end{center}%
  \end{figure}%
}

%
% Comando para incluir imágenes en formato apaisado (sin marco).
\newcommand{\figuraApaisadaSinMarco}[5]{%
  \begin{figure}%
    \begin{center}%
    \includegraphics[angle=90,height=#1\textheight,#5]{#2}%
    \caption{#3}%
    \label{#4}%
    \end{center}%
  \end{figure}%
}
% Para las tablas
\newcommand{\otoprule}{\midrule [\heavyrulewidth]}
%
% Nuevo comando para tablas pequeñas (menos de una página).
\newcommand{\tablaSmall}[5]{%
 \begin{table}
  \begin{center}
   \rowcolors {2}{gray!35}{}
   \begin{tabular}{#2}
    \toprule
    #4
    \otoprule
    #5
    \bottomrule
   \end{tabular}
   \caption{#1}
   \label{tabla:#3}
  \end{center}
 \end{table}
}

%
% Nuevo comando para tablas pequeñas (menos de una página).
\newcommand{\tablaSmallSinColores}[5]{%
 \begin{table}[H]
  \begin{center}
   \begin{tabular}{#2}
    \toprule
    #4
    \otoprule
    #5
    \bottomrule
   \end{tabular}
   \caption{#1}
   \label{tabla:#3}
  \end{center}
 \end{table}
}

\newcommand{\tablaApaisadaSmall}[5]{%
\begin{landscape}
  \begin{table}
   \begin{center}
    \rowcolors {2}{gray!35}{}
    \begin{tabular}{#2}
     \toprule
     #4
     \otoprule
     #5
     \bottomrule
    \end{tabular}
    \caption{#1}
    \label{tabla:#3}
   \end{center}
  \end{table}
\end{landscape}
}

%
% Nuevo comando para tablas grandes con cabecera y filas alternas coloreadas en gris.
\newcommand{\tabla}[6]{%
  \begin{center}
    \tablefirsthead{
      \toprule
      #5
      \otoprule
    }
    \tablehead{
      \multicolumn{#3}{l}{\small\sl continúa desde la página anterior}\\
      \toprule
      #5
      \otoprule
    }
    \tabletail{
      \hline
      \multicolumn{#3}{r}{\small\sl continúa en la página siguiente}\\
    }
    \tablelasttail{
      \hline
    }
    \bottomcaption{#1}
    \rowcolors {2}{gray!35}{}
    \begin{xtabular}{#2}
      #6
      \bottomrule
    \end{xtabular}
    \label{tabla:#4}
  \end{center}
}

%
% Nuevo comando para tablas grandes con cabecera.
\newcommand{\tablaSinColores}[6]{%
  \begin{center}
    \tablefirsthead{
      \toprule
      #5
      \otoprule
    }
    \tablehead{
      \multicolumn{#3}{l}{\small\sl continúa desde la página anterior}\\
      \toprule
      #5
      \otoprule
    }
    \tabletail{
      \hline
      \multicolumn{#3}{r}{\small\sl continúa en la página siguiente}\\
    }
    \tablelasttail{
      \hline
    }
    \bottomcaption{#1}
    \begin{xtabular}{#2}
      #6
      \bottomrule
    \end{xtabular}
    \label{tabla:#4}
  \end{center}
}

%
% Nuevo comando para tablas grandes sin cabecera.
\newcommand{\tablaSinCabecera}[5]{%
  \begin{center}
    \tablefirsthead{
      \toprule
    }
    \tablehead{
      \multicolumn{#3}{l}{\small\sl continúa desde la página anterior}\\
      \hline
    }
    \tabletail{
      \hline
      \multicolumn{#3}{r}{\small\sl continúa en la página siguiente}\\
    }
    \tablelasttail{
      \hline
    }
    \bottomcaption{#1}
  \begin{xtabular}{#2}
    #5
   \bottomrule
  \end{xtabular}
  \label{tabla:#4}
  \end{center}
}



\definecolor{cgoLight}{HTML}{EEEEEE}
\definecolor{cgoExtralight}{HTML}{FFFFFF}

%
% Nuevo comando para tablas grandes sin cabecera.
\newcommand{\tablaSinCabeceraConBandas}[5]{%
  \begin{center}
    \tablefirsthead{
      \toprule
    }
    \tablehead{
      \multicolumn{#3}{l}{\small\sl continúa desde la página anterior}\\
      \hline
    }
    \tabletail{
      \hline
      \multicolumn{#3}{r}{\small\sl continúa en la página siguiente}\\
    }
    \tablelasttail{
      \hline
    }
    \bottomcaption{#1}
    \rowcolors[]{1}{cgoExtralight}{cgoLight}

  \begin{xtabular}{#2}
    #5
   \bottomrule
  \end{xtabular}
  \label{tabla:#4}
  \end{center}
}



\graphicspath{ {./img/} }

% Capítulos
\chapterstyle{bianchi}
\newcommand{\capitulo}[2]{
	\setcounter{chapter}{#1}
	\setcounter{section}{0}
	\setcounter{figure}{0}
	\setcounter{table}{0}
	\chapter*{\thechapter.\enskip #2}
	\addcontentsline{toc}{chapter}{\thechapter.\enskip #2}
	\markboth{#2}{#2}
}

% Apéndices
\renewcommand{\appendixname}{Apéndice}
\renewcommand*\cftappendixname{\appendixname}

\newcommand{\apendice}[1]{
	%\renewcommand{\thechapter}{A}
	\chapter{#1}
}

\renewcommand*\cftappendixname{\appendixname\ }

% Formato de portada
\makeatletter
\usepackage{xcolor}
\newcommand{\tutor}[1]{\def\@tutor{#1}}
\newcommand{\course}[1]{\def\@course{#1}}
\definecolor{cpardoBox}{HTML}{E6E6FF}
\def\maketitle{
  \null
  \thispagestyle{empty}
  % Cabecera ----------------
\noindent\includegraphics[width=\textwidth]{cabecera}\vspace{1cm}%
  \vfill
  % Título proyecto y escudo informática ----------------
  \colorbox{cpardoBox}{%
    \begin{minipage}{.8\textwidth}
      \vspace{.5cm}\Large
      \begin{center}
      \textbf{TFG del Grado en Ingeniería Informática}\vspace{.6cm}\\
      \textbf{\LARGE\@title{}}
      \end{center}
      \vspace{.2cm}
    \end{minipage}

  }%
  \hfill\begin{minipage}{.20\textwidth}
    \includegraphics[width=\textwidth]{escudoInfor}
  \end{minipage}
  \vfill
  % Datos de alumno, curso y tutores ------------------
  \begin{center}%
  {%
    \noindent\LARGE
    Presentado por \@author{}\\ 
    en Universidad de Burgos --- \@date{}\\
    Tutor: \@tutor{}\\
  }%
  \end{center}%
  \null
  \cleardoublepage
  }
\makeatother

\newcommand{\nombre}{David Cezar Toderas} %%% cambio de comando

% Datos de portada
\title{
	 Detección de \textit{malware} mediante técnicas de Inteligencia Artificial}
\author{\nombre}
\tutor{Álvar Arnaiz González}
\date{\today}

\begin{document}

\maketitle


\newpage\null\thispagestyle{empty}\newpage


%%%%%%%%%%%%%%%%%%%%%%%%%%%%%%%%%%%%%%%%%%%%%%%%%%%%%%%%%%%%%%%%%%%%%%%%%%%%%%%%%%%%%%%%
\thispagestyle{empty}


\noindent\includegraphics[width=\textwidth]{cabecera}\vspace{1cm}

\noindent D. Álvar Arnaiz González, profesor del departamento de Ingeniería Informática, Área de Lenguajes y Sistemas Informáticos.

\noindent Expone:

\noindent Que el alumno D. \nombre, con DNI X8300130M, ha realizado el Trabajo final de Grado en Ingeniería Informática titulado título de TFG. 

\noindent Y que dicho trabajo ha sido realizado por el alumno bajo la dirección del que suscribe, en virtud de lo cual se autoriza su presentación y defensa.

\begin{center} %\large
En Burgos, {\large \today}
\end{center}

\vfill\vfill\vfill

% Author and supervisor
%\begin{minipage}{0.45\textwidth}
%\begin{flushleft} %\large
%Vº. Bº. del Tutor:\\[2cm]
%D. nombre tutor
%\end{flushleft}
%\end{minipage}
%\hfill
%\begin{minipage}{0.45\textwidth}
%\begin{flushleft} %\large
%Vº. Bº. del co-tutor:\\[2cm]
%D. nombre co-tutor
%\end{flushleft}
%\end{minipage}
%\hfill
%
%\vfill

% para casos con solo un tutor comentar lo anterior
% y descomentar lo siguiente
\begin{center}
	Vº. Bº. del Tutor:\\[2cm]
	Álvar Arnaiz González	
\end{center}


\newpage\null\thispagestyle{empty}

\textbf{Agradecimientos}

Este trabajo se ha construido sobre dos pilares: el apoyo incondicional de mi familia, que ha sido mi ancla, y la infinita paciencia y confianza de mi tutor, que ha sido mi guía. A mis amigos, gracias por entender que mi ausencia no era olvido, sino la dedicación que este proyecto exigía, y por seguir ahí con el mismo ánimo de siempre. A todos vosotros, por ser el motor, el refugio y la razón por la que cada hora invertida en estas páginas ha merecido la pena.

\newpage

\frontmatter

% Abstract en castellano
\renewcommand*\abstractname{Resumen}
\begin{abstract}
	
El crecimiento exponencial del sistema operativo Android lo ha convertido en el principal objetivo del \textit{software} malicioso. Este trabajo aborda el desarrollo de un sistema de detección de \textit{malware} para Android basado en inteligencia artificial y análisis estático, un enfoque que permite evaluar la seguridad de las aplicaciones sin necesidad de ejecutarlas. El proyecto se desarrolló en dos fases: una primera de prototipado para validar la viabilidad del enfoque, y una segunda de ingeniería en la que se construyó un \textit{pipeline} completo, incluyendo la creación de un \textit{dataset} propio con 20\,000 aplicaciones y el desarrollo de un modelo de red neuronal a medida.

Los resultados demuestran la alta eficacia del sistema, con modelos capaces de alcanzar un \textit{recall} y un \textit{$F_1$-Score} superiores al 98\%. La conclusión más relevante del trabajo es el papel fundamental del \textit{embedder} de la red neuronal, un componente que transforma los datos brutos en una representación numérica de alta calidad. Este \textit{embedder} no solo permite que la red neuronal funcione, sino que, al ser utilizado como preprocesador para algoritmos de aprendizaje automático clásicos (como XGBoost o Regresión Logística), potencia su rendimiento a niveles excepcionales, demostrando el poder de las arquitecturas híbridas.

\end{abstract}

\renewcommand*\abstractname{Descriptores}
\begin{abstract}
detección de \textit{malware}, android, análisis estático, inteligencia artificial, redes neuronales, aprendizaje profundo, \textit{embeddings}, interpretabilidad de modelos, SHAP, ciberseguridad.
\end{abstract}

\clearpage

% Abstract en inglés
\renewcommand*\abstractname{Abstract}
\begin{abstract}

The exponential growth of the Android operating system has made it the primary target for malicious software. This project addresses the development of a malware detection system for Android based on artificial intelligence and static analysis, an approach that allows for the assessment of application security without the need for execution. The project was carried out in two phases: an initial prototyping phase to validate the approach's feasibility, and a second engineering phase in which a complete pipeline was built, including the creation of a custom dataset with 20,000 applications and the development of a tailored neural network model.

The results demonstrate the high effectiveness of the system, with models capable of achieving recall and F1-Score metrics exceeding 98\%. The most significant conclusion of this work is the fundamental role of the neural network's embedder, a component that transforms raw data into a high-quality numerical representation. This embedder not only enables the neural network to function but also, when used as a preprocessor for classical machine learning algorithms (such as XGBoost or Logistic Regression), boosts their performance to exceptional levels, showcasing the power of hybrid architectures.

\end{abstract}

\renewcommand*\abstractname{Keywords}
\begin{abstract}
malware detection, android, static analysis, artificial inteligence, neural network, deep learning, embeddings, model interpretability, SHAP, cibersecurity.
\end{abstract}

\clearpage

% Indices
\tableofcontents

\clearpage

\listoffigures

\clearpage

\listoftables
\clearpage

\mainmatter
\capitulo{1}{Introducción}


% TODO: Introducción al documento

% Visión general del problema, Android es una plataforma mundialmente usada que compone la gran mayoría del mercado de los smartphones hoy en día. (añadir cifra y cita)

% Explicación del proyecto y lo que se ha realizado, las distitnas fases de este.

% Breve explicación de los resultados

Intro

\section{Estructura de la memoria}

El trabajo se divide en dos partes principales, una memoria (este documento) con la explicación general de lo realizado y, un documento adicional de anexos que contiene diferentes apéndices relativos al proyecto que indagan en cuestiones más específicas acerca del mismo. La estructura de ambos documentos puede verse reflejada a continuación:

\subsection{Memoria}

La memoria (este documento) sigue la siguiente estructura:

\begin{enumerate}
	\item \textbf{Introducción:} Presenta el proyecto, su contexto y motivación, y describe la estructura del documento.
	
	\item \textbf{Objetivos del proyecto:} Detalla las metas específicas que el proyecto busca alcanzar, tanto a nivel de funcionalidades como de requisitos técnicos.
	
	\item \textbf{Conceptos teóricos:} Establece el marco teórico necesario para poder comprender el documento.
	
	\item \textbf{Técnicas y herramientas:} Describe el conjunto de librerías y otras tecnologías que se han empleado para desarrollar el proyecto.
	
	\item \textbf{Aspectos relevantes del desarrollo del proyecto:} Narra el proceso realizado a lo largo del trabajo, explicando las fases, las decisiones tomadas y cómo se llegó a la solución.
	
	\item \textbf{Trabajos relacionados:} Análisis del estado del arte relacionado con el proyecto.
	
	\item \textbf{Conclusiones y líneas de trabajo futuras:} Reflexiones e ideas finales acerca del proyecto y aspectos a mejorar en un futuro.
\end{enumerate}

\subsection{Anexos}

El documento de anexos tiene los siguientes apéndices:

\begin{enumerate}[label=\Alph*]
	\item \textbf{Plan de proyecto software:} Contiene la planificación del proyecto y los estudios de viabilidad de este.
	 
	\item \textbf{Especificación de requisitos:} Expone los requisitos funcionales y no funcionales del sistema, catálogo de requisitos y especificaciones adicionales.
	
	\item \textbf{Especificación de diseño:} Describe el diseño de datos y arquitectónico del proyecto.
	
	\item \textbf{Documentación técnica de programación:} Guía destinada a los desarrolladores que explica la estructura del código, las decisiones de implementación y los pasos para configurar el entorno de desarrollo.
	
	\item \textbf{Documentación de usuario:} Manual de instrucciones dirigido al usuario final, que explica cómo utilizar la aplicación y sus funcionalidades.
\end{enumerate}

\subsection{Recursos adicionales}

Además de la memoria y los anexos, se proporciona el enlace a la aplicación creada y al repositorio del proyecto.

\begin{enumerate}
	\item \textbf{Aplicación web:} 
	
	\item \textbf{Repositorio del proyecto:} 
\end{enumerate}

\capitulo{2}{Objetivos del proyecto}

Este apartado explica de forma precisa y concisa cuáles son los objetivos que se persiguen con la realización del proyecto. Se puede distinguir entre los objetivos marcados por los requisitos del software a construir y los objetivos de carácter técnico que plantea a la hora de llevar a la práctica el proyecto.

\section{Objetivos funcionales}
Los objetivos funcionales definen las capacidades y acciones que el sistema final debe ofrecer al usuario. Son, en esencia, las características con las que se podrá interactuar directamente.

\begin{itemize}
	\item \textbf{Clasificación de aplicaciones Android:} El sistema debe ser capaz de recibir un archivo APK y analizarlo para determinar si es malicioso o benigno, devolviendo un resultado claro.
	
	\item \textbf{Desarrollo de una interfaz web de demostración:} Crear una aplicación web simple e intuitiva que permita a un usuario subir un archivo APK y obtener la predicción del modelo de forma visual.
	
	\item \textbf{Interpretabilidad del modelo:} Ser capaz de explicar el porqué de las decisiones del modelo y reconocer cuales son los factores que más afectan a este y entender cómo realiza las predicciones que estima.
\end{itemize}

\section{Objetivos técnicos}
Los objetivos técnicos se refieren a las metas de implementación, arquitectura y proceso necesarias para construir el sistema y asegurar su calidad y funcionalidad interna.

\begin{itemize}
	\item \textbf{Investigación del estado del arte:} Realizar un análisis de la literatura científica existente para comprender las técnicas actuales de detección de \textit{malware} con IA y posicionar el proyecto.
	
	\item \textbf{Creación de un dataset propio:} Diseñar y construir un conjunto de datos a medida, extrayendo características útiles de una cantidad suficiente de ficheros APK.
	
	\item \textbf{Implementación de un \textit{pipeline} de extracción de características:} Desarrollar un proceso completo y replicable que transforme los datos brutos de una APK en el formato numérico que el modelo necesita para su análisis.
	
	\item \textbf{Diseño y entrenamiento del modelo de red neuronal:} Construir un modelo de \textit{deep learning} con PyTorch, entrenarlo con el \textit{dataset} propio y optimizar su rendimiento para que sea capaz de clasificar \textit{malware} correctamente.
	
	\item \textbf{Ajuste de hiperparámetros:} Utilizar herramientas de búsqueda automatizada, como Optuna, para encontrar la configuración de hiperparámetros que ofrezca el mejor equilibrio entre rendimiento y tiempo de entrenamiento.
	
	\item \textbf{Cuantización del modelo:} Intentar reducir el tamaño del modelo todo lo posible para optimizar su rendimiento y facilitar su despliegue en otras plataformas.
	
	\item \textbf{Análisis comparativo de modelos:} Entrenar clasificadores clásicos de aprendizaje automático (como KNN, SVM, Linear Regression, Random Forest o XGBoost) usando las características procesadas por la red neuronal y comparar sus resultados.
\end{itemize}

\section{Objetivos personales}
Estos objetivos describen las competencias y conocimientos que se esperan adquirir a nivel personal durante el desarrollo del proyecto.

\begin{itemize}
	\item \textbf{Aplicar conocimientos de IA en un proyecto real:} Llevar la teoría a la práctica, diseñando, entrenando y evaluando un modelo de red neuronal desde cero con Python.
	
	\item \textbf{Adquirir experiencia en investigación académica:} Aprender a buscar, leer y sintetizar documentación científica, tanto para fundamentar las decisiones técnicas del proyecto, como para mejorar mi compresión a la hora de leer estudios de cualquier tipo.
	
	\item \textbf{Profundizar en IA y ciberseguridad:} Ganar un mejor entendimiento acerca de estos dos grandes campos que hoy en día dominan el mundo de la informática. Comprender a su vez mejor cómo funciona el \textit{malware} y las técnicas que pueden usarse para detectarlo y combatirlo.
	
	\item \textbf{Obtener un modelo útil y con un buen rendimiento:} Ser capaz de entrenar un modelo que pueda competir con otros. El objetivo principal en este caso es obtener un modelo que sea capaz de acertar en la mayoría de los casos y que dé predicciones bastante fiables.
\end{itemize}

\capitulo{3}{Conceptos teóricos}

En aquellos proyectos que necesiten para su comprensión y desarrollo de unos conceptos teóricos de una determinada materia o de un determinado dominio de conocimiento, debe existir un apartado que sintetice dichos conceptos.

Algunos conceptos teóricos de \LaTeX{} \footnote{Créditos a los proyectos de Álvaro López Cantero: Configurador de Presupuestos y Roberto Izquierdo Amo: PLQuiz}.

\section{Secciones}

Las secciones se incluyen con el comando section.

\subsection{Subsecciones}

Además de secciones tenemos subsecciones.

\subsubsection{Subsubsecciones}

Y subsecciones. 


\section{Referencias}

Las referencias se incluyen en el texto usando cite~\cite{wiki:latex}. Para citar webs, artículos o libros~\cite{koza92}, si se desean citar más de uno en el mismo lugar~\cite{bortolot2005, koza92}.


\section{Imágenes}

Se pueden incluir imágenes con los comandos standard de \LaTeX, pero esta plantilla dispone de comandos propios como por ejemplo el siguiente:

\imagen{escudoInfor}{Autómata para una expresión vacía}{.5}



\section{Listas de items}

Existen tres posibilidades:

\begin{itemize}
	\item primer item.
	\item segundo item.
\end{itemize}

\begin{enumerate}
	\item primer item.
	\item segundo item.
\end{enumerate}

\begin{description}
	\item[Primer item] más información sobre el primer item.
	\item[Segundo item] más información sobre el segundo item.
\end{description}
	
\begin{itemize}
\item 
\end{itemize}

\section{Tablas}

Igualmente se pueden usar los comandos específicos de \LaTeX o bien usar alguno de los comandos de la plantilla.

\tablaSmall{Herramientas y tecnologías utilizadas en cada parte del proyecto}{l c c c c}{herramientasportipodeuso}
{ \multicolumn{1}{l}{Herramientas} & App AngularJS & API REST & BD & Memoria \\}{ 
HTML5 & X & & &\\
CSS3 & X & & &\\
BOOTSTRAP & X & & &\\
JavaScript & X & & &\\
AngularJS & X & & &\\
Bower & X & & &\\
PHP & & X & &\\
Karma + Jasmine & X & & &\\
Slim framework & & X & &\\
Idiorm & & X & &\\
Composer & & X & &\\
JSON & X & X & &\\
PhpStorm & X & X & &\\
MySQL & & & X &\\
PhpMyAdmin & & & X &\\
Git + BitBucket & X & X & X & X\\
Mik\TeX{} & & & & X\\
\TeX{}Maker & & & & X\\
Astah & & & & X\\
Balsamiq Mockups & X & & &\\
VersionOne & X & X & X & X\\
} 

\capitulo{4}{Técnicas y herramientas}

% TODO: Los links a las distintas aplicaciones deberían ir en el pie de página o en la bibliografía; lo mismo con la plantilla de latex

En este apartado se comentan y analizan las diferentes herramientas y técnicas usadas en la realización de este proyecto.

\section{Lenguajes y añadidos}

En esta sección se describe el lenguaje de programación principal y los conceptos y herramientas de entorno que han sido fundamentales para la organización y el desarrollo del proyecto.

\subsection{Python}

% https://www.python.org/

Python es un lenguaje de programación alto nivel, interpretado y multipropósito, conocido por su sintaxis clara y sencilla muy similar al inglés. Su filosofía de diseño se centra en torno a la facilidad de uso, lo que lo ha convertido en uno de los lenguajes más populares del mundo para todo tipo de personas que desean tanto adentrarse en el mundo de la programación, como aquellas que ya tienen cierta experiencia en el campo, especialmente si se trata de tareas de automatización, de ciencia de datos o incluso de IA. Especialmente en este último grupo, prácticamente todo el ecosistema de IA moderno se desarrolla sobre Python. Todo esto se debe, principalmente a su enorme comunidad, su facilidad para crear prototipos rápidos y, sobre todo, su vasto repertorio de librerías especializadas y optimizadas para el cálculo científico y el aprendizaje automático.

\subsection{Entorno virtual}

% https://peps.python.org/pep-0405/

Un entorno virtual es un concepto de programación, especialmente popular en Python, que consiste en crear un directorio aislado que contiene una instalación propia de Python y todas las dependencias específicas de un proyecto dado. Funciona como una <<burbuja>> que separa los paquetes de un proyecto de los de otros proyectos y de la instalación global del sistema. Esto resuelve el clásico problema de tener diferentes proyectos que requieren versiones distintas de las mismas librerías.

La principal ventaja de usar entornos virtuales es que son fáciles de reproducir y de limpiar. Permiten saber con exactitud qué dependencias necesita un proyecto y evitan <<contaminar>> el sistema con paquetes que solo se usan para una tarea concreta. Cuando un proyecto finaliza, basta con borrar la carpeta del entorno para eliminar todo rastro de este, sin dejar dependencias huérfanas. Sin embargo, su gestión tradicional con herramientas como \texttt{venv} y \texttt{pip} puede ser manual y algo tediosa, ya que uno debe encargarse de activarlos, desactivarlos y registrar las dependencias explícitamente.

\subsection{Poetry}

% https://python-poetry.org/

Poetry es una herramienta de gestión de dependencias y empaquetado para Python. A diferencia de los métodos tradicionales que combinan herramientas como \texttt{pip}, \texttt{requirements.txt} y \texttt{venv}, Poetry integra todas estas funcionalidades bajo una única interfaz de comandos y un archivo de configuración llamado \texttt{pyproject.toml}. Dicho archivo puede ser modificado para declarar las dependencias del proyecto, características del paquete final, versiones y otro tipo de configuraciones adicionales para la generación de distribuciones y reglas de resolución de paquetes. Poetry se encargará posteriormente de generar a partir de dicho archivo de configuración un fichero llamado \texttt{poetry.lock}, el cual contiene todas las versiones exactas de cada paquete y sus dependencias correspondientes, correctamente resueltas y guardadas en un formato que permitirá posteriormente reproducir el entorno de forma simple y directa.

En este caso, Poetry no solo es útil para proyectos grandes sino que es extremadamente cómodo para cualquier proyecto de Python que necesite de un par de dependencias, puesto que permite crear entornos virtuales con solo un par de comandos, permitiendo así no modificar la instalación global de Python en el sistema y permitiendo una gestión más compleja de ciertas dependencias en función de las prestaciones del equipo que uno posee o su sistema operativo. Por ejemplo, en este proyecto, una de las librerías usadas es PyTorch, la cual cuenta con soporte para aceleración por GPU de sus operaciones, pero, este soporte viene incluido en un paquete diferente al del paquete que solo permite el uso de la CPU. Usando Poetry es posible crear una configuración dinámica que sea capaz de adaptarse al equipo de cada uno, descargando o no la versión con o sin aceleración por hardware en función de si el equipo es compatible con ello.

\subsection{Conda}

% https://anaconda.org/anaconda/conda

Conda es un sistema de gestión de paquetes y entornos de código abierto y multiplataforma. Aunque puede gestionar cualquier lenguaje, es extremadamente popular en la comunidad de Python, especialmente para todo lo relacionado con ciencia de datos e IA. Su distribución principal, Anaconda, viene con un enorme conjunto de librerías científicas preinstaladas, facilitando su uso y dando acceso a todo lo necesario para empezar a realizar pruebas y desarrollar proyectos desde el momento en el que se instala. Por otro lado, existe también una versión mínima, Miniconda, que solo incluye el gestor de entornos y permite al usuario instalar únicamente lo que necesita, siendo una opción más ligera y la preferida por muchos simplemente porque, a pesar de requerir un poco más de tiempo y conocimiento en preparar el entorno de trabajo, uno puede instalar específicamente solo lo que necesita.

La gran ventaja de Conda es que gestiona no solo paquetes de Python, sino también dependencias que no son de Python (como librerías de C o compiladores), lo que simplifica mucho la instalación de paquetes complejos como PyTorch. Sin embargo, una de sus desventajas es que utiliza sus propios repositorios de paquetes, que a veces no están tan actualizados como el repositorio oficial de Python (PyPI). Además, compartir y replicar entornos de forma exacta puede ser menos directo que con herramientas más modernas como Poetry.

\subsection{Jupyter Notebook}
% Entorno de desarrollo de Python interactivo que se equipara a una especie de cuaderno interactivo. En este "cuaderno" el usuario es capaz de mezclar tanto código python como texto excrito con Markdown para realizar explicaciones acerca del código o explicar ciertos añadidos a este. Es muy utlizado sobretodo en la comunidad científica y cuando se quiere realizar prototipado de modelos principalemente porque es muy simple, todo lo que haces se guarda de forma temporal en memoria y puedes reeditar tu código o añadir más código sin tener que reejecutar todo este desde cero, ahorrando el hecho de por ejemplo, reentrenar un modelo entero si lo único que querías cambiar era cómo se ve una gráfica de resultados o, tener que cargar de vuelta en memoria todo un dataset grande para comprobar si la forma en la que lo procesas es correcta. Es muy útil para este tipo de tareas donde reejetuar ciertas operaciones es costoso y uno quiere realizar pruebas rápidas con ciertos componentes o funcionalidades.
% En este caso se ha usado principalemente para realizar experimentos de todo tipo a lo largo del proyecto, desde probar a cómo filtrar los datos, a la construcción del modelo, mejora de este, pruebas de descarga de apks para la creación del dataset propio, pruebas con la extracción de características, la contrucción del mismo dataset, la representación de los resultados e inclso pruebas de análsis e interpretabilidad. Es una herramietna muy util que ahoorra mucho tiempo en este tipo de actividades de prototipado. Aunque tiene sus errores pero sirve muy bien para su propósito.

% https://jupyter.org/

Un Jupyter Notebook es un entorno de desarrollo interactivo basado en una aplicación web que permite crear y compartir documentos que contienen código mixto, ecuaciones, visualizaciones y texto. Funciona como un <<cuaderno de laboratorio digital>>, donde se puede escribir y ejecutar código en bloques o "celdas" de forma independiente. Esto es extremadamente útil para el prototipado y el análisis exploratorio de datos.

Su principal ventaja es la interactividad. Permite ejecutar una celda costosa (como cargar un gran dataset o entrenar un modelo) una sola vez, y luego seguir trabajando en otras celdas (como visualizar resultados o probar transformaciones) sin tener que volver a ejecutar todo el script desde el principio. Esto ahorra una cantidad de tiempo enorme, especialmente en este tipo de tareas de experimentación.

Jupyter ha sido una herramienta indispensable a lo largo de todo este proyecto, especialmente en las fases de prototipado y análisis. Se utilizó para experimentar con la extracción de características, para depurar el proceso de creación del \textit{dataset}, para entrenar las primeras versiones del modelo, para visualizar los resultados con Matplotlib e incluso para ejecutar los análisis de interpretabilidad con SHAP. Fue el <<banco de pruebas>> donde se validaron la mayoría de las ideas antes de integrarlas en el código final del proyecto.

\section{Librerías}

En esta sección se describen las librerías de Python más importantes que se han utilizado a lo largo de todo el proyecto.

\subsection{PyTorch vs Keras}

% https://pytorch.org/
% https://keras.io/

PyTorch y Keras representan dos de los \textit{frameworks} de \textit{deep learning} más conocidos y utilizados en la actualidad. Si bien ambos facilitan la construcción de redes neuronales, cada uno toma una filosofía de diseño distinta, lo cual los hace más adecuados para diferentes tipos de proyectos y usuarios. Keras funciona como una interfaz de alto nivel, diseñada para la simplicidad y el desarrollo rápido, mientras que PyTorch opera a un nivel más bajo, ofreciendo un control granular y una mayor flexibilidad a la hora de desarrollar modelos personalizados y específicos.

Keras se caracteriza principalmente por su facilidad de uso, debido a que permite construir y entrenar modelos estándar con muy pocas líneas de código. Su API es bastante intuitiva y abstrae gran parte de la complejidad subyacente, lo que lo convierte en una opción excelente para principiantes y para la creación rápida de prototipos y para entornos de producción donde la estandarización es clave. Sin embargo, esta simplicidad conlleva una menor flexibilidad, ya que realizar modificaciones sustanciales en la arquitectura o en el ciclo de entrenamiento de los modelos puede volverse complejo y poco intuitivo debido a no estar diseñado para ello.

Por el contrario, PyTorch proporciona un set de herramientas mucho más versátil, pensado para la investigación y para proyectos que requieren de arquitecturas personalizadas. En general, permite crear modelos de forma mucho más granular debido a que se basa en proporcionar al usuario con un conjunto de módulos y funciones que pueden ser instanciadas juntas para formar un modelo complejo y personalizado, especializado en la tarea que se desee. A su vez, permite definir de manera más concreta el set de datos a usar, el preprocesado de esos datos y el cómo se entrena el modelo a partir de ellos.

En concreto, para este proyecto, PyTorch fue la elección más lógica. La arquitectura del modelo de clasificación de APKs no es convencional; se procesan un montón de cadenas de caracteres y vectores, los cuales requieren de un preprocesado específico y de la implementación de un \textit{embedder} personalizado para poder usar posteriormente dichas características junto con una lógica y control sobre el proceso de entrenamiento bastante concreto.

% TODO: Preguntar acercade "soporte" como meter una nueva linea

\tablaSmall{Comparativa entre PyTorch y Keras}{p{.25\textwidth} | p{.33\textwidth} | p{.33\textwidth}}{comparativapytorchkeras}
{\textbf{Características} & \textbf{Pytoch} & \textbf{Keras} \\}
{
\textbf{Flexibilidad} & Alta, permite crear redes neuronales personalizadas y complejas. & Menos flexible, enfocado en redes estándar. \\
\textbf{Facilidad de uso} & Requiere más código y tiene una curva de aprendizaje más pronunciada. & Fácil de usar, ideal para desarrolladores principiantes. \\
\textbf{Personalización} & Excelente para redes personalizadas y modelos avanzados. & Limitada, más orientada a redes convencionales. \\
\textbf{Comunidad y soporte} & Muy popular en la investigación académica y proyectos avanzados. & Amplio uso en la industria por su simplicidad. \\
\textbf{Uso principal} & Investigación, redes neuronales complejas. & Desarrollo rápido de modelos estándar. \\
}

\subsection{NumPy, Pandas y Matplotlib}

% https://numpy.org/
% https://pandas.pydata.org/
% https://matplotlib.org/stable/

NumPy, Pandas y Matplotlib constituyen la trilogía fundamental de librerías sobre las que se edifica gran parte del ecosistema de paquetes científicos y de data science en Python.

\begin{itemize}
	\item \textbf{NumPy (\textit{Numerical Python}):} Es la librería base para la computación numérica en Python. Proporciona el concepto de ndarray, una estructura de datos para la creación de arrays N-dimensionales eficientes, y un vasto conjunto de funciones matemáticas para operar sobre ellos. Una de las mayores ventajas de NumPy es su velocidad y eficiencia tanto en tiempo, como en espacio al trabajar con grandes sets de datos. Este rendimiento se debe principalmente a que muchas de sus operaciones están implementadas en C y aprovechan la vectorización, permitiendo ejecutar operaciones complejas en arrays completos sin necesidad de bucles explícitos en Python.
	
	\item \textbf{Pandas:} Construida sobre NumPy, Pandas introduce estructuras de datos de alto nivel, principalmente el DataFrame, una tabla bidimensional heterogénea e indexada, pensada para manejar datos tabulares y series temporales. Facilita enormemente tareas como la lectura y escritura de datos, la limpieza, el filtrado, la agregación y la transformación, siendo una herramienta muy útil estándar para el preprocesamiento de datos.

	\item \textbf{Matplotlib:} Es la librería de visualización de datos por excelencia en Python. Permite elegir entre un gran repertorio de gráficos comunes como pueden ser los gráficos de barras o histogramas a la creación de gráficos personalizados, facilitando el trabajo de representar datos complejos y hacerlos agradables a la vista.
	
\end{itemize}

Estas tres librerías han sido de gran ayuda en diferentes etapas del proyecto. Pandas ha sido la herramienta principal para estructurar las características extraídas de los archivos APK en un DataFrame limpio y manejable, facilitando todo el preprocesamiento. NumPy ha sido utilizado de forma subyacente por Pandas y PyTorch, y directamente para realizar operaciones numéricas eficientes sobre los datos ya procesados antes de introducirlos en el modelo. Finalmente, Matplotlib ha sido usada para la evaluación del modelo, permitiendo visualizar diferentes aspectos del entrenamiento del modelo y la comparación de este con otros de una manera más visual.

\subsection{scikit-learn}
% Librería para entrenamiento de modelos clásicos de ml, proporciona una gran base de ejemplos y todo tipo de herramietnas y utilidades para trabajar con datos, entrenar y evaluar modelos de diferentes tipos, la mayor ventaja que posee es que todo los modelos y la gran mayoría de sus técnicas caen dentro de una interfaz común la cual es muy simple de entender y usar.
% Usada para entrenar los modelos clásicos de ML (SVM, KNN, RandomForest y LogisticReg, a su vez se puede mencionar xgboost, es un paquete a parte pero está basado en scikit-learn). Además a sido muy útil en el proyecto por muchas de sus funcioes de utilidad a la hora de manejar grnades sets de datos, principalemente para el muestreo de forma justa y reproducible del dataet.

% https://scikit-learn.org/stable/

scikit-learn es la librería de referencia para el aprendizaje automático clásico en Python. Proporciona un conjunto enorme y bien documentado de herramientas para prácticamente cualquier tarea de \textit{machine learning}, incluyendo algoritmos de clasificación, regresión, \textit{clustering}, reducción de dimensionalidad y utilidades para el preprocesamiento de datos y la evaluación de modelos. Su mayor ventaja es su API unificada y consistente que facilita en gran medida el trabajo de entrenar modelos y trabajar con ellos: todos sus objetos (modelos, transformadores, etc.) comparten una interfaz común (\texttt{.fit()}, \texttt{.predict()}, \texttt{.transform()}), lo que hace que sea muy fácil de aprender y usar.

A lo largo del proyecto, scikit-learn ha sido de gran ayuda a la hora de comparar la red neuronal con otros modelos clásicos. Se utilizó para entrenar y evaluar dicho modelos (SVM, $k$-NN, RandomForest, Regresión Logística) con los que poder comprar posteriormente la red. Además, sus funciones de utilidad, como las de división de datos (\texttt{train\_test\_split}) y cálculo de métricas (\texttt{accuracy\_score()}, \texttt{precision\_score()}, \texttt{roc\_curve()}, ...), fueron de gran ayuda para poder obtener cifras acerca del rendimiento de los distintos clasificadores.

\subsection{Optuna}
% Tuning del modelo, búsqueda inteligente de los mejores hiperparámetros para este. Explicar levemente el método que usa y decir que es muy utilizado hoy en día con redes neuroanel especificamente porque es muy cómodo a la hora de cómo bucsa la mejor combinación de parámetros. No es como un simple GridSearch o RandomSearch de scikit o algo similar, sino que construye una especie de minimodelo que aprende cómo se comporta el modelo de forma iterativa a medida que lo entrena, esto permite a Optuna cortar de forma agreisa todos aquellos caminos (decisiones) de hiperparámetros que afectan negativamente al modelo y quedarse con el más óptimo, centrando las pruebas cada vez más en obtener lo mejor combinación posble. Es una especie de búsqeuda guiada.
% Se usa principalemtne porque es simple y provee una interfaz muy sencilla para trabajr con ello, solo es necesario crear un pequeño wrapper al modelo con los rangos de hiperpaámetros a probar y el resto se encarga optuna de realizarlo. Aunque si que es verdad  que es necesario tener cuidado proque el modelo y sus processo de entrenamiento tienen que estar bien hecho y tener en cuetna que se deberán de poder correr de forma secuencial. Esto implica que si, la librería nos da mucha flexibilidad pero también implica que tenemos que procurar que ciertas cosas como la limpieza de memoria y el error handling de ciertos tests sea hecho por noosotros.

Optuna~\cite{optuna_2019} es un \textit{framework} de optimización de hiperparámetros moderno que sirve de sustituto a métodos de búsqueda tradicionales basados en fuerza bruta por una estrategia de optimización secuencial e informada. Su diseño <<\textit{define-by-run}>> lo hace extremadamente flexible y le permite integrarse de forma simple con variedad de \textit{frameworks} de aprendizaje automático como es el caso de PyTorch.

Por defecto, el algoritmo que Optuna utiliza internamente es el \textit{Tree-structured Parzen Estimator}~\cite{watanabe2023tree} (TPE). Este método funciona de una manera bastante ingeniosa: en lugar de modelar directamente la probabilidad de que unos hiperparámetros den un buen resultado, se modelan dos distribuciones de probabilidad distintas. La primera, $l(x)$, representa la distribución de los hiperparámetros que han dado lugar a buenos resultados (pérdida baja), y la segunda, $g(x)$, la de aquellos que han dado lugar a malos resultados. En cada nuevo paso, Optuna busca un conjunto de hiperparámetros que sea muy probable bajo el modelo <<bueno>> y muy poco probable bajo el modelo <<malo>>, maximizando así la probabilidad de encontrar una configuración cada vez mejor. Este enfoque dirigido es mucho más eficiente que una búsqueda aleatoria y más óptimo que una búsqueda exhaustiva porque poda una gran cantidad de casos que no son necesarios probar si se conoce que parte de sus hiperparámetros no generan buenos resultados.

Dado que el rendimiento de una red neuronal depende enormemente de sus hiperparámetros (tasa de aprendizaje, número de capas, etc.), se eligió Optuna para automatizar y optimizar este proceso de búsqueda. Su sencilla integración con PyTorch permitió definir un espacio de búsqueda y dejar que Optuna explorara de forma eficiente cientos de combinaciones, encontrando una configuración casi óptima en mucho menos tiempo de lo que habría llevado un enfoque manual o mediante una búsqueda exhaustiva.

\subsection{SHAP (\textit{SHapley Additive exPlanations})}
% Liberría que permite implementar un método de análisis de modelo llamado, Shapley addivide values, es una librería uy popualr puesto que esta técnica permite tanto realizar análisis locales a una predicción concreta de los modelos como de forma global a toda una población de datos. Explicar el concepto detrás de las Shap values (citar parer), mostrar las ecucaiones del problema para un ejemplo simple (el ejemplo clásico del que parte, teoría de juegos, repartir un monto de dinero entre unos participanetes de manera justa en base a su contribución al concurso), explicar luego como esto se relaciona con el ml y porque se puede hacer casi directamente la asociación de este problema a asociar el valor de la predicción final a cada una de las diferentes características de entrada en base a su contribución. Explicar cambio de las ecuaciones frente al ejemplo clásico (citar paper acerca de la adaptación a ml) y decir que la gran ventaja que esta librería aporta es el hecho de aproximaar muy bien y rápidamente el cálculo de las Shapely values. Es importante también aclaarar que las Shap values no implican directamente causalidad en la vida real (citar paper que descubrió esto) sino que solo sirven apra determinar en qué se fija el modelo, esto no se puede extrapolar a que eso sea debido a que en la vida real esa característica a su vez sea así de importante o sea justo la más decisva, simplemente evalua el cómo funciona el modelo. A su vez, es verdad que es un buen inicio para intentar buscar correlaciones pero se tiene que estudiar el tema con cuidado. 
%Es importante destacar que esta herramietna sirve principalemente para permitir interpretar las predicciones de los distintos modelos que tenemos y para determinar en qué se fija cada uno de ellos apra realizar sus predicciones.

% https://github.com/slundberg/shap

SHAP (\textit{SHapley Additive exPlanations}) es una técnica de interpretabilidad de modelos que responde a una pregunta fundamental: si un modelo ha tomado una decisión, ¿cuánto ha contribuido cada una de las características de entrada a ese resultado final? Para ello, se basa en un concepto de la teoría de juegos cooperativos llamado los valores de Shapley~\cite{lundberg2017unified}. La analogía es simple: si un equipo gana un premio, ¿cómo se reparte el dinero de forma justa entre sus miembros, teniendo en cuenta la aportación de cada uno? SHAP permite resolver exactamente este problema, pero equiparando a los <<jugadores>> y al <<juego>> con las características de un modelo y su predicción respectivamente.

El algoritmo funciona de una manera teóricamente muy elegante. Para calcular la contribución de una característica (por ejemplo, <<la aplicación solicita acceso a los contactos>>), SHAP considera todas las combinaciones posibles del resto de características (las <<coaliciones>>). Luego, mide cuánto cambia la predicción del modelo cuando se añade esa característica a cada una de esas coaliciones. Al promediar este <<cambio marginal>> a través de todas las coaliciones posibles, se obtiene la contribución justa y única de esa característica a la predicción final. El resultado es una explicación aditiva: la suma de las contribuciones de todas las características nos da exactamente la diferencia entre la predicción concreta y la predicción media del modelo.

SHAP se utilizó para analizar e interpretar las predicciones de los distintos modelos entrenados. Permitiendo visualizar qué características (permisos, llamadas a la API del sistema, etc.) eran las más influyentes para cada modelo a la hora de clasificar una aplicación como \textit{malware}. Además, debido a su naturaleza aditiva e individual, SHAP puede ser utilizado tanto para el análisis local de una predicción concreta, como para obtener una explicación general del rendimiento del modelo y del porque de sus decisiones, convirtiéndolo en una herramienta muy útil y versátil para analizar los distintos modelos que se entrenaron.

\subsection{UMAP (\textit{Uniform Manifold Approximation and Projection})}

UMAP (\textit{Uniform Manifold Approximation and Projection})~\cite{mcinnes2018umap} es un algoritmo moderno de reducción de dimensionalidad basado en principios matemáticos de la topología y el aprendizaje de variedades (\textit{manifold learning}). La idea fundamental de UMAP es que la mayoría de los conjuntos de datos de alta dimensión en realidad viven en una <<variedad>> o superficie de una dimensionalidad mucho más baja que está <<curvada>> dentro de ese espacio. La misión de UMAP es encontrar esa variedad y <<desplegarla>> en un espacio de pocas dimensiones (como una hoja de papel 2D) perdiendo la menor cantidad de información estructural posible, facilitando así su visualización e interpretabilidad.

Para lograrlo, su algoritmo se basa en dos etapas. Primero, construye una representación topológica de los datos en el espacio de alta dimensión. Para ello, crea un grafo ponderado donde los pesos de las aristas se calculan asumiendo que la distancia entre los puntos varía localmente a lo largo de la variedad. En la segunda etapa, busca una <<incrustación>> (\textit{embedding}) de baja dimensión de ese grafo, es decir, una disposición de los puntos en 2D o 3D, que tenga una estructura topológica lo más parecida posible. Para medir esta <<similitud>> entre el grafo de alta y baja dimensión, utiliza una función de coste basada en la entropía cruzada, que optimiza para encontrar la mejor proyección posible.

Se utilizó UMAP frente a otras técnicas como t-SNE por su equilibrio entre rendimiento y calidad de la visualización. Es computacionalmente más rápido y, en general, hace un mejor trabajo preservando la estructura global de los datos. En este caso, esto era ideal para visualizar el espacio de \textit{embeddings} y confirmar que el modelo que se entrena no solo agrupaba correctamente las muestras a nivel local, sino que la separación global entre la clase <<\textit{malware}>> y la <<\textit{benigna}>> era clara y coherente.

\subsection{Streamlit}

% https://streamlit.io/

Streamlit es un \textit{framework} de código abierto para Python, diseñado específicamente para la creación y el despliegue rápido de aplicaciones web interactivas para proyectos de \textit{data science} y aprendizaje automático. Su filosofía se basa en la simplicidad radical, permitiendo transformar \textit{scripts} convencionales con código de procesamiento de datos, modelos de IA o simples funciones aisladas en aplicaciones web funcionales con un esfuerzo y conocimiento de desarrollo web mínimos.

A diferencia de \textit{frameworks} web más tradicionales y complejos como Django o Flask, que exigen la gestión de rutas, una organización de los ficheros del proyecto específica, sintaxis inusual, plantillas HTML y lógica de servidor, Streamlit permite construir una interfaz de usuario directamente desde un script normal de Python. Con comandos sencillos, se pueden añadir elementos interactivos como botones, deslizadores, gráficos y, fundamentalmente para este caso, campos para la subida de archivos. Esto acelera drásticamente el ciclo de desarrollo puesto que uno puede simplemente centrarse en obtener un modelo o set de funcionalidades que son correctas y dan buenos resultados sin preocuparse mucho de cómo se llevará luego esto a una interfaz gráfica o aplicación de escritorio / web puesto que la conversión es muy sencilla en la mayoría de casos.

La finalidad de la aplicación web en este proyecto es ofrecer una demostración tangible y una especie de \textit{demo} o entorno de prueba para que cualquiera pudiera probar el clasificador de \textit{malware} de forma sencilla. El objetivo era crear una interfaz simple donde un usuario pudiera subir un archivo \texttt{.apk} y recibir una predicción de manera inmediata. Streamlit fue la herramienta perfecta para esta tarea, ya que permitió desarrollar esta funcionalidad en cuestión de horas en lugar de días y el resultado es más que suficiente para el alcance deseado.

\subsection{Androguard y otros analizadores}

% https://github.com/androguard/androguard
% https://github.com/MobSF/Mobile-Security-Framework-MobSF
% https://github.com/skylot/jadx

Androguard es una potente herramienta de código abierto y un paquete de Python, diseñada específicamente para el análisis estático y la ingeniería inversa de aplicaciones de Android (archivos APK). Su función principal es el proceso de diseccionar un archivo APK para extraer información detallada sobre su estructura y contenido sin necesidad de ejecutar la aplicación. Una de las mayores ventajas de Androguard es el hecho de que permite además realizar todo este proceso de extracción de características de forma automatizada puesto que, expone una API de Python bastante simple e intuitiva que proporciona acceso a todos los datos que se pueden obtener de analizar las APKs.

A través de Androguard, es posible acceder a componentes como el \textit{manifest} de la aplicación (AndroidManifest.xml) para analizar permisos y componentes declarados, desensamblar el código Dalvik (DEX) para inspeccionar las clases y los métodos e, incluso, extraer recursos como cadenas de texto o certificados. Aunque existen otras herramientas de análisis como MobSF (que ofrece un entorno más automatizado y visual) o Jadx (un popular descompilador), Androguard destaca por su granularidad, simplicidad de uso y su naturaleza como librería de Python.

Una de las piedras fundamentales de este proyecto es la posibilidad de extraer características de forma estática de APKs y entrenar un modelo con ellas capaz de discernir entre aplicaciones benignas y malignas. La razón principal de su eso es el hecho de funcionar nativamente en Python y permitir la automatización de la creación del \textit{dataset} de entrenamiento del modelo. A su vez, es el componente que permite analizar muestras nuevas, obteniendo los datos que el modelo espera recibir para realizar una predicción.

\section{Otras herramientas}

Finalmente, en esta sección se detallan otras herramientas de software que, aunque no están directamente relacionadas con la inteligencia artificial, han sido de gran ayuda para la gestión del código, el desarrollo y la documentación del proyecto.

\subsection{Docker}

% https://www.docker.com/

Docker es una plataforma de código abierto que permite automatizar el despliegue, la ejecución, la distribución y la gestión de aplicaciones mediante el uso de la ''containerización''. Esta herramienta permite empaquetar una aplicación y todas sus dependencias como bibliotecas, herramientas del sistema, código, comandos de ejecución, configuraciones y todo lo que uno pueda necesitar para ejecutar su aplicación en una unidad independiente y aislada llamada contenedor.

El principal problema que Docker intenta resolver es el clásico ''en mi equipo funciona'', donde una aplicación se ejecuta correctamente en el entorno de un desarrollador pero falla en otro debido a diferencias en la configuración del sistema operativo o en las versiones de las dependencias, además de posibles errores en la ejecución del proyecto y diversas posibles causas. Un contenedor soluciona esto debido a que, uno puede simplemente distribuir un único archivo que describe el proceso de creación del contenedor llamado \textit{Dockerfile} (o la misma imagen del contenedor ya creado) garantizando que cualquiera podrá ejecutar el proyecto exactamente de la misma manera en cualquier máquina que soporte Docker, desde un portátil local hasta un servidor en la nube. Esto asegura la consistencia, la reproducibilidad y simplifica enormemente los flujos de trabajo desarrollo.

El entorno de este proyecto es complejo, requiriendo versiones específicas de Python y múltiples librerías (PyTorch, Androguard, etc.). Para asegurar que la aplicación web y el modelo pudieran ser ejecutados por cualquier persona sin un tedioso proceso de configuración manual, se utilizó Docker para empaquetar todo en un contenedor. Esto garantiza que el proyecto funcione de la misma manera en cualquier sistema, simplificando enormemente su despliegue y asegurando que los resultados fueran reproducibles.

\subsection{GitHub}
% Usado para hostear el repositorio de código del proyecto. Explicar levemetne que es un repositorio de códdigo también.

% https://github.com/
% https://github.com/dtx1007/tfg_24_25

GitHub es una plataforma de desarrollo colaborativo basada en la web que utiliza el famosos sistema de control de versiones llamado Git. Permite a los desarrolladores alojar y gestionar sus repositorios de código, que son esencialmente carpetas de un proyecto con un historial completo de todos los cambios realizados. Además de alojar el código, GitHub ofrece herramientas para el seguimiento de errores, la revisión de código y la gestión de proyectos, facilitando el trabajo en equipo.

Todo el código fuente del proyecto ha sido alojado en un repositorio de GitHub. Esto sirvió principalmente para tener un espacio compartido desde el cual poder trabajar desde diferentes equipos manteniendo la coherencia de los archivos. Además, sirvió como copia de seguridad por si acaso algo salía mal y era necesario recuperar los archivos de versiones anteriores.

\subsection{Visual Studio Code (VSCode)}
% Comentar lo que es un IDE y porqué es muy útil a la hora de desarrolar y simplificar el trabajo del desarrollador. Es un IDE modular, tiene un montón de extensiones y es bastante liviano y rápdio lo cual lo permite en una herramienta con una base muy buena y que además es extensible, haciendo que sea de gran ayuda en la gran mayoría de casos. Mencionar a Jetbrains como una alternativa muy buena y en concreto a Pytorch que es su editor más orientado hacia la programación en Python.
% IDE usado para editar el código 

% https://code.visualstudio.com/
% https://www.jetbrains.com/es-es/pycharm/

Visual Studio Code (VSCode) es un editor de código fuente liviano y simple muy utililzado en el mundo del desarrollo gracias a su diseño mínimo, velocidad su carácter modular. El editor en sí es bastante completo y sirve perfectamente para su propósito pero, lo que lo eleva al siguiente nivel es su enorme ecosistema de extensiones que permiten adaptarlo para trabajar con prácticamente cualquier lenguaje o tecnología. Una alternativa muy popular en el mundo de Python es la suite de JetBrains, especialmente su IDE PyCharm.

VSCode fue el editor de código principal utilizado para el desarrollo de todo el proyecto. Su flexibilidad, su excelente soporte para Python y sus extensiones para Docker y Jupyter lo convirtieron en la herramienta ideal para gestionar las diferentes facetas del trabajo desde un único lugar.

\subsection{TeXstudio}

% https://www.texstudio.org/

TeXstudio es un Entorno de Desarrollo Integrado (IDE) diseñado específicamente para la creación de documentos con \LaTeX. \LaTeX\ es un sistema de composición de textos de alta calidad, muy utilizado en el mundo académico y científico para la escritura de artículos, tesis y libros, ya que maneja de forma excepcional las fórmulas matemáticas, las referencias cruzadas y la maquetación profesional. TeXstudio facilita el trabajo con \LaTeX\ al proporcionar un visor de PDF integrado, autocompletado de comandos, resaltado de sintaxis y otras herramientas que agilizan y simplifican el proceso de escritura.

Toda la documentación de este proyecto, incluyendo esta memoria y sus anexos, ha sido escrita en \LaTeX utilizando la plantilla proporcionada por la universidad~\cite{UBU2024ubutfgm}. TeXstudio fue la herramienta elegida para esta tarea, ya que su entorno integrado hizo mucho más cómodo y eficiente el proceso de redactarla.

\capitulo{5}{Aspectos relevantes del desarrollo del proyecto}

Este apartado pretende recoger los aspectos más interesantes del desarrollo del proyecto, comentados por los autores del mismo.
Debe incluir desde la exposición del ciclo de vida utilizado, hasta los detalles de mayor relevancia de las fases de análisis, diseño e implementación.
Se busca que no sea una mera operación de copiar y pegar diagramas y extractos del código fuente, sino que realmente se justifiquen los caminos de solución que se han tomado, especialmente aquellos que no sean triviales.
Puede ser el lugar más adecuado para documentar los aspectos más interesantes del diseño y de la implementación, con un mayor hincapié en aspectos tales como el tipo de arquitectura elegido, los índices de las tablas de la base de datos, normalización y desnormalización, distribución en ficheros3, reglas de negocio dentro de las bases de datos (EDVHV GH GDWRV DFWLYDV), aspectos de desarrollo relacionados con el WWW...
Este apartado, debe convertirse en el resumen de la experiencia práctica del proyecto, y por sí mismo justifica que la memoria se convierta en un documento útil, fuente de referencia para los autores, los tutores y futuros alumnos.

\capitulo{6}{Trabajos relacionados}

Este apartado sería parecido a un estado del arte de una tesis o tesina. En un trabajo final grado no parece obligada su presencia, aunque se puede dejar a juicio del tutor el incluir un pequeño resumen comentado de los trabajos y proyectos ya realizados en el campo del proyecto en curso. 

\capitulo{7}{Conclusiones y Líneas de trabajo futuras}

Este capítulo final tiene como objetivo recapitular los hallazgos y resultados obtenidos a lo largo del desarrollo del proyecto, ofreciendo una reflexión crítica sobre el cumplimiento de los objetivos y las lecciones aprendidas. Asimismo, se propondrán diversas líneas de trabajo futuras que podrían servir como continuación de esta investigación, explorando nuevas vías para mejorar, optimizar y ampliar la solución desarrollada.

\section{Conclusiones}

El desarrollo de este proyecto ha culminado en la creación de un sistema completo y funcional para la detección de \textit{malware} en Android mediante análisis estático e inteligencia artificial. A lo largo de este recorrido, se han obtenido una serie de conclusiones técnicas y conceptuales de gran relevancia.

La primera y más evidente conclusión es que \textbf{el enfoque propuesto es altamente eficaz}. Todos los modelos entrenados, tanto la red neuronal como los clasificadores clásicos, han demostrado una capacidad de detección muy buena, alcanzando métricas de \textit{recall}, precisión y AUC superiores al 98\% en la mayoría de los casos. Esto confirma de manera rotunda la hipótesis central del proyecto: es perfectamente viable construir un clasificador de alto rendimiento basándose únicamente en las características estáticas extraídas de un archivo APK, sin necesidad de recurrir a técnicas de análisis dinámico, que son más lentas y costosas. El objetivo de obtener un modelo útil y con buen rendimiento se ha cumplido con creces.

Sin embargo, la conclusión más importante y reveladora de este trabajo es el \textbf{papel fundamental del \textit{embedder} como motor de ingeniería de características}. Si bien la red neuronal por sí sola es un clasificador muy competente, su verdadero valor radica en su capacidad para aprender una representación numérica densa y significativa a partir de datos de entrada complejos y heterogéneos. La visualización del espacio de \textit{embeddings} mediante UMAP demostró de forma sorprendente que este componente era capaz de separar casi perfectamente las aplicaciones benignas de las maliciosas en un espacio de alta dimensionalidad. Este preprocesamiento inteligente es la razón por la cual modelos mucho más simples, como la Regresión Logística o XGBoost, pudieron alcanzar un rendimiento tan espectacular. El proyecto demuestra en la práctica un principio fundamental del aprendizaje automático: la calidad de los datos y de sus representaciones es, a menudo, más importante que la complejidad del modelo clasificador final. Esto valida la estrategia híbrida como un enfoque extremadamente poderoso.

En tercer lugar, el análisis de interpretabilidad mediante SHAP ha permitido \textbf{abrir la <<caja negra>> de los modelos y validar su proceso de razonamiento}. Los resultados mostraron que todos los clasificadores, de manera consistente, basaban sus decisiones en las características más lógicas: los permisos solicitados, las llamadas a la API y los componentes declarados en la aplicación (actividades, servicios y receptores). Esto no solo aumenta la confianza en las predicciones, sino que también cumple el objetivo de crear un modelo interpretable, demostrando que no estaba aprendiendo de artefactos o sesgos extraños en los datos, sino de indicadores de comportamiento reales.

Desde una perspectiva de ingeniería, otra conclusión importante es que la \textbf{optimización y el diseño del \textit{pipeline} de datos son tan críticos como el propio modelo}. La transición del prototipo al modelo final reveló graves cuellos de botella relacionados con el manejo de la memoria y el preprocesamiento de datos, cuya solución requirió bastante trabajo de optimización. Esto subraya la importancia de entender no solo los algoritmos, sino también las herramientas de software (Pandas, NumPy, PyTorch) a un nivel más bajo para construir sistemas eficientes y escalables.

Finalmente, la exploración de la \textbf{cuantización del modelo arrojó resultados limitados para esta arquitectura específica}. Se concluyó que las técnicas de cuantización post-entrenamiento no son efectivas cuando la mayor parte del peso de un modelo reside en las capas de \textit{embedding}. Esto supone una conclusión técnica valiosa: para llevar un modelo de estas características a entornos con recursos limitados, como un dispositivo móvil, se requerirían enfoques más avanzados, como el entrenamiento consciente de la cuantización y con ello, el posible rediseño de la arquitectura.

En resumen, se han cumplido todos los objetivos principales del proyecto. Se ha realizado una investigación del estado del arte, se ha construido un \textit{dataset} propio con un \textit{pipeline} replicable, se ha entrenado y optimizado una red neuronal de alto rendimiento y se ha realizado un análisis comparativo e interpretativo exhaustivo, culminando en una aplicación web funcional. A nivel personal, el proyecto ha permitido adquirir una profunda experiencia práctica en todo el ciclo de vida de un proyecto de IA, desde la investigación y la gestión de datos hasta el despliegue.

\section{Líneas de trabajo futuras}

Todo proyecto de investigación y desarrollo abre nuevas puertas y plantea nuevas preguntas. Este trabajo no es una excepción. A continuación, se detallan varias líneas de trabajo futuras que podrían dar continuidad y expandir las ideas y soluciones aquí presentadas.

La primera y más obvia línea de futuro sería el \textbf{despliegue y la evaluación del modelo en un dispositivo de Android real}. Aunque la aplicación web es una excelente herramienta de demostración, el objetivo final de un sistema de este tipo sería poder analizar las aplicaciones directamente en el teléfono del usuario. Esto presenta importantes desafíos de optimización, puesto que, como se concluyó en el análisis de cuantización, el modelo actual es demasiado grande (cerca de 1 GB). Una línea de investigación futura muy interesante sería explorar técnicas avanzadas para reducir su tamaño, como el \textit{pruning} (poda de conexiones neuronales), el entrenamiento consciente de la cuantización (\textit{Quantization-Aware Training}) o incluso el rediseño del modelo utilizando arquitecturas específicamente pensadas para dispositivos móviles, como las inspiradas en MobileNet, adaptándolas para procesar las características estáticas en lugar de imágenes.

Una segunda vía de expansión natural sería la de \textbf{escalar el modelo y el \textit{dataset}}. Durante la búsqueda de hiperparámetros, se encontró una configuración de modelo mucho más grande y con mayor capacidad teórica de aprendizaje, pero que no ofrecía una mejora de rendimiento significativa con el \textit{dataset} actual de 20\,000 aplicaciones. Una línea de trabajo futura consistiría en ampliar de forma masiva el \textit{dataset}, utilizando el \textit{pipeline} ya creado para procesar muchas más APKs de AndroZoo. Con un volumen de datos mucho mayor, sería posible que el modelo más grande sí fuera capaz de aprender patrones más sutiles y complejos, superando el rendimiento del modelo actual y de los clasificadores clásicos.

En tercer lugar, se podría \textbf{enriquecer el conjunto de características extraídas}. Aunque el análisis estático ha demostrado ser muy potente, la incorporación de otros tipos de información podría mejorar aún más la precisión. Se podría explorar la inclusión de características basadas en grafos, como los grafos de llamadas a la API o los grafos de flujo de control, y utilizar modelos de Redes Neuronales de Grafos (GNN) para analizarlos, siguiendo la línea de trabajos como el de MAPAS~\cite{kim2022mapas}. Otra opción aún más ambiciosa sería evolucionar hacia un <<modelo híbrido>>, incorporando un número limitado de características dinámicas obtenidas de la ejecución de la aplicación en un \textit{sandbox}, para complementar la visión estática con información acerca del comportamiento real de las aplicaciones.

Finalmente, la \textbf{aplicación web de demostración podría ser mejorada y expandida}. Se podrían añadir funcionalidades como un historial de análisis persistente para cada usuario, al igual que un sistema de usuarios y de gestión que permitiera añadir y modificar los modelos que estarían disponibles en la aplicación. Otro posible enfoque sería reestructurar la aplicación y usar un \textit{framework} web más completo como puede ser Django. Esto convertiría la herramienta de una simple demostración a una plataforma de análisis educativo y de diagnóstico mucho más completa.




\bibliographystyle{plain}
\bibliography{bibliografia}

\end{document}
